\clearpage
\begin{trapbox}
\textbf{\textcolor{CTrap}{误区一：外层是偶函数，整体一定偶。——错}}
\begin{center}\begin{tikzpicture}[font=\small,>=stealth,line join=round]
\draw[gray!40,rounded corners=3pt] (0,0) rectangle (14.2,8.3);\node at (7.1,7.700000000000001) {外层平方为偶，整体却偏离了纵轴};
\begin{scope}[shift={(5.1,2.068181818181818)},x=1.3333333333333333cm,y=0.8363636363636363cm]\draw[->,gray] (-3,0)--(1.5,0) node[right]{$x$};\draw[->,gray] (0,-0.5)--(0,5) node[above]{$y$};\draw[blue,thick,solid,domain=-2.8:1.2,samples=180] plot(\x,{(\x+1)^2});\fill[orange!85!black] (-1,0) circle(2.2pt);\fill[orange!85!black] (1,4) circle(2.2pt);\end{scope}\node[anchor=west,align=left,text width=5.0cm] at (8.4,4.0) {$F=(x+1)^2$\\[10pt]内层 $g=x+1$ 非奇非偶\\[10pt]$F(1)=4,\ F(-1)=0$\\[10pt]对称轴 $x=-1$，不是 $y$ 轴};
\node[text width=12.6cm,align=center] at (7.1,.65) {外层只会把互为相反数的输入算成相同结果；这里的2和0并非相反数。};
\end{tikzpicture}\end{center}


反例 $f(u)=u^2$，$g(x)=x+1$，得到 $F(x)=(x+1)^2$。虽然外面平方能把 $u$ 与 $-u$ 算成相同数，但\textbf{里面给到外面的两个数未必正好相反}：$g(1)=2$、$g(-1)=0$，平方后分别为 4 和 0。不能只看外层。

\medskip

\end{trapbox}
\clearpage
\begin{trapbox}
\textbf{\textcolor{CTrap}{误区二：只要里面是奇函数，整体一定奇。——错}}
\begin{center}\begin{tikzpicture}[font=\small,>=stealth,line join=round]
\draw[gray!40,rounded corners=3pt] (0,0) rectangle (14.2,8.3);\node at (7.1,7.700000000000001) {例3：内层虽反号，外层平方把两路合并};
\begin{scope}[shift={(4.1,2.161111111111111)},x=1.875cm,y=0.5111111111111111cm]\draw[->,gray] (-1.6,0)--(1.6,0) node[right]{$x$};\draw[->,gray] (0,-1)--(0,8) node[above]{$y$};\draw[blue,thick,solid,domain=-1.3:1.3,samples=180] plot(\x,{4*(\x)^2});\fill[orange!85!black] (-1,4) circle(2.2pt);\fill[orange!85!black] (1,4) circle(2.2pt);\end{scope}\node[anchor=west,align=left,text width=5.0cm] at (8.4,4.0) {$F=(2x)^2$\\[10pt]$g(1)=2,\ g(-1)=-2$\\[10pt]平方以后都变成 $4$\\[10pt]内奇外偶，整体偶};
\node[text width=12.6cm,align=center] at (7.1,.65) {不能把“内层为奇”误记成“整体必奇”；还需要检查外层。};
\end{tikzpicture}\end{center}


例如 $g(x)=x$ 是奇函数，$f(u)=u^2$ 是偶函数，整体 $F(x)=x^2$ 为偶函数。\textbf{记住口诀后半句“同外”，还要看外层。}

\medskip

\end{trapbox}
\clearpage
\begin{trapbox}
\textbf{\textcolor{CTrap}{误区三：整体是偶函数，内层必是偶函数。——错}}
\begin{center}\begin{tikzpicture}[font=\small,>=stealth,line join=round]
\draw[gray!40,rounded corners=3pt] (0,0) rectangle (14.2,8.3);\node at (7.1,7.700000000000001) {例3：内层虽反号，外层平方把两路合并};
\begin{scope}[shift={(4.1,2.161111111111111)},x=1.875cm,y=0.5111111111111111cm]\draw[->,gray] (-1.6,0)--(1.6,0) node[right]{$x$};\draw[->,gray] (0,-1)--(0,8) node[above]{$y$};\draw[blue,thick,solid,domain=-1.3:1.3,samples=180] plot(\x,{4*(\x)^2});\fill[orange!85!black] (-1,4) circle(2.2pt);\fill[orange!85!black] (1,4) circle(2.2pt);\end{scope}\node[anchor=west,align=left,text width=5.0cm] at (8.4,4.0) {$F=(2x)^2$\\[10pt]$g(1)=2,\ g(-1)=-2$\\[10pt]平方以后都变成 $4$\\[10pt]内奇外偶，整体偶};
\node[text width=12.6cm,align=center] at (7.1,.65) {不能把“内层为奇”误记成“整体必奇”；还需要检查外层。};
\end{tikzpicture}\end{center}


例如 $F(x)=x^2=f(g(x))$，其中 $g(x)=x$ 是奇函数，$f(u)=u^2$ 是偶函数。整体偶，不代表内层偶。甚至内外均不具备奇偶性时，整体仍可能出现对称。

\medskip

\end{trapbox}
\clearpage
\begin{trapbox}
\textbf{\textcolor{CTrap}{误区四：外层既非奇也非偶，就一定无法判断。——错}}
\begin{center}\begin{tikzpicture}[font=\small,>=stealth,line join=round]
\draw[gray!40,rounded corners=3pt] (0,0) rectangle (14.2,8.3);\node at (7.1,7.700000000000001) {外层没有奇偶性，内偶仍能保证整体偶};
\begin{scope}[shift={(4.1,2.068181818181818)},x=1.875cm,y=0.8363636363636363cm]\draw[->,gray] (-1.6,0)--(1.6,0) node[right]{$x$};\draw[->,gray] (0,-0.5)--(0,5) node[above]{$y$};\draw[blue,thick,solid,domain=-1.4:1.4,samples=180] plot(\x,{pow(2,(\x)^2)});\fill[orange!85!black] (-1,2) circle(2.2pt);\fill[orange!85!black] (1,2) circle(2.2pt);\end{scope}\node[anchor=west,align=left,text width=5.0cm] at (8.4,4.0) {$F=2^{x^2}$\\[10pt]$x$ 与 $-x$ 先得到同一指数\\[10pt]外层 $2^u$ 非奇非偶\\[10pt]最后函数值仍然相同};
\node[text width=12.6cm,align=center] at (7.1,.65) {两路先合并为同一输入，才是“内偶则偶”的决定因素。};
\end{tikzpicture}\end{center}


如果内层是偶函数，根本不需要外层的奇偶性。例如 $F(x)=2^{x^2}$ 为偶函数。另外，只要肯直接算 $F(-x)$，很多口诀管不了的题仍能判断。

\medskip

\end{trapbox}
\clearpage
\begin{trapbox}
\textbf{\textcolor{CTrap}{误区五：忘了检查定义域。——错}}
\begin{center}\begin{tikzpicture}[font=\small,>=stealth,line join=round]
\draw[gray!40,rounded corners=3pt] (0,0) rectangle (14.2,8.3);\node at (7.1,7.700000000000001) {定义域只向一侧延伸，不能判奇偶};
\begin{scope}[shift={(3.9125,2.4166666666666665)},x=0.9375cm,y=1.5333333333333332cm]\draw[->,gray] (-3,0)--(3.4,0) node[right]{$x$};\draw[->,gray] (0,-0.5)--(0,2.5) node[above]{$y$};\draw[blue,thick,solid,domain=-1:3,samples=180] plot(\x,{sqrt(\x+1)});\fill[orange!85!black] (-1,0) circle(2.2pt);\end{scope}\node[anchor=west,align=left,text width=5.0cm] at (8.4,4.0) {$F=\sqrt{x+1}$\\[10pt]$D=[-1,+\infty)$\\[10pt]$2$ 能代入，$-2$ 不能\\[10pt]既非奇也非偶};
\node[text width=12.6cm,align=center] at (7.1,.65) {先看整体定义域；不能只因表达式里有根号就猜测对称性。};
\end{tikzpicture}\end{center}


看 $F(x)=\sqrt{x+1}$，它只允许 $x\geqslant-1$，不关于 $0$ 左右对称，所以既不是奇函数也不是偶函数。不能只凭公式外形猜对称。\textbf{另一个对照：} $\sqrt{x^2-1}$ 的定义域是 $(-\infty,-1]\cup[1,+\infty)$，虽然中间缺了一段，却左右对称，仍可判为偶函数。

\medskip

\end{trapbox}
\clearpage
\begin{trapbox}
\textbf{\textcolor{CTrap}{误区六：用某一对数代入成功，就宣布函数是奇或偶。——错}}
\begin{center}\begin{tikzpicture}[font=\small,>=stealth,line join=round]
\draw[gray!40,rounded corners=3pt] (0,0) rectangle (14.2,8.3);\node at (7.1,7.700000000000001) {一个零点满足关系，不能保证整张图像对称};
\begin{scope}[shift={(4.1,2.6184210526315788)},x=1.875cm,y=0.8070175438596491cm]\draw[->,gray] (-1.6,0)--(1.6,0) node[right]{$x$};\draw[->,gray] (0,-1.2)--(0,4.5) node[above]{$y$};\draw[blue,thick,solid,domain=-1.3:1.3,samples=180] plot(\x,{(\x)^2+(\x)^3});\fill[orange!85!black] (-1,0) circle(2.2pt);\fill[orange!85!black] (0,0) circle(2.2pt);\fill[orange!85!black] (1,2) circle(2.2pt);\end{scope}\node[anchor=west,align=left,text width=5.0cm] at (8.4,4.0) {$H=x^2+x^3$\\[10pt]虽然 $H(0)=H(-0)=0$\\[10pt]但 $H(1)=2,\ H(-1)=0$\\[10pt]既非奇也非偶};
\node[text width=12.6cm,align=center] at (7.1,.65) {一个输入可用于寻找反例，不能替代对每个允许输入的证明。};
\end{tikzpicture}\end{center}


例如 $H(x)=x^2+x^3$，取 $x=0$ 时 $H(0)=H(-0)=0$，但这不能证明它是偶函数。判定奇偶性要求\textbf{每一个允许的 $x$} 都符合相应关系。代入一个数只能帮助发现反例，不能代替完整证明。

\medskip

\end{trapbox}
\clearpage
\begin{trapbox}
\textbf{\textcolor{CTrap}{误区七：把“内奇同外”理解成任何情况下都完全一样。——不严谨}}
\begin{center}\begin{tikzpicture}[font=\small,>=stealth,line join=round]
\draw[gray!40,rounded corners=3pt] (0,0) rectangle (14.2,8.3);\node at (7.1,7.700000000000001) {同样得到偶函数，外层的性质可以完全不同};
\begin{scope}[shift={(4.1,2.068181818181818)},x=2.0cm,y=0.41818181818181815cm]\draw[->,gray] (-1.5,0)--(1.5,0) node[right]{$x$};\draw[->,gray] (0,-1)--(0,10) node[above]{$y$};\draw[blue,thick,solid,domain=-1.2:1.2,samples=180] plot(\x,{(\x)^2+1});\draw[violet,thick,dashed,domain=-1.2:1.2,samples=180] plot(\x,{((\x)^2+1)^2+(\x)^2+1});\end{scope}\node[anchor=west,align=left,text width=5.0cm] at (8.4,4.0) {内层都为 $g=x^2+1$\\[10pt]蓝：外层 $f(u)=u$（奇）\\[10pt]紫：$f(u)=u^2+u$（非奇非偶）\\[10pt]两个整体都是偶函数};
\node[text width=12.6cm,align=center] at (7.1,.65) {两条实际曲线说明：仅从整体为偶，无法反推外层的奇偶性。};
\end{tikzpicture}\end{center}


如果外层已知是奇函数，整体一定奇；外层已知是偶函数，整体一定偶。但是如果外层既非奇也非偶，整体也可能因为只用到了外层某一部分取值而出现奇偶性。\textbf{口诀只说明这些条件足够保证结论，不允许随便倒过来推。}

\medskip
\textbf{\textcolor{CTrap}{最后检查清单}}
\begin{enumerate}[leftmargin=2.2em,itemsep=0.2em]
\item $x$ 与 $-x$ 都能代入整个式子吗？
\item 先算 $g(-x)$，它等于 $g(x)$，还是 $-g(x)$？
\item 如需看外层，$f(-u)$ 等于 $f(u)$ 还是 $-f(u)$？
\item 如果都不符合，直接写出 $F(-x)$ 进行比较，不要强套口诀。
\end{enumerate}
\end{trapbox}
